\clearpage
\fancyhead[L]{\small VON ZACH. TABULAE ABERRATIONIS ETC.}
\hypertarget{native-494-lower}{}\label{native-494-lower}
\begin{center}
\small Göttingische gelehrte Anzeigen. Stück 2. S. 18..24. 2. Januar 1808.\\[3pt]
\rule{.72\linewidth}{.3pt}
\end{center}

\emph{Tabulae speciales aberrationis et nutationis in ascensionem rectam et in declinationem ad supputandas stellarum fixarum positiones, sive apparentes, sive veras, una cum insigniorum 494 stellarum zodiacalium catalogo novo in specula astronomica Ernestina ad initium anni \textup{1800} constructo, cum aliis tabulis eo spectantibus, auctore} \textsc{Franc. Lib. Bar. de Zach}. Vol. I. 1806. 208 und \textsc{clvi} S. Vol. II. 1807. 508 S. Großoktav.

Das Erscheinen dieses wichtigen Werks, das die wertvollen Früchte langjähriger Arbeit enthält, wird allen praktischen Astronomen höchst willkommen sein. Sie haben es schon seit mehreren Jahren erwartet. Es behandelt vor allem die kleinen Veränderungen der Positionen der Fixsterne durch Präzession, Aberration und Nu%
\SourceLeaf{495}%
tation. Diese Positionen bilden bekanntlich die Grundlage der meisten astronomischen Beobachtungen. Die Berechnung jener kleinen Veränderungen gehört daher zu den täglichen Aufgaben des praktischen Astronomen. Um diese Astronomen hat sich der berühmte Verfasser also mit einem Werk sehr verdient gemacht, das die flexibelsten und brauchbarsten Formeln für jene Veränderungen aufstellt, durch äußerst sorgfältige Erörterungen anhand der neuesten und bewährtesten Untersuchungen und Beobachtungen die zuverlässigsten Zahlenwerte der zugrunde liegenden Größen ermittelt und die Arbeit durch zweckmäßige und bequeme Hilfstafeln so weit wie möglich verkürzt. Mit einem Wort, es bietet alles, was man sich zu diesem Gegenstand wünschen kann. So wertvoll dies alles ist, möchte der Rezensent einen anderen beträchtlichen Teil des Werks noch höher einschätzen, obwohl er im Titel nicht einmal erwähnt ist. Es ist das Verzeichnis von 1830 Tierkreissternen, ein Denkmal, dauerhafter als Erz, für zwei Sternwarten, die eine Zierde Deutschlands waren und jetzt leider beide ungenutzt dastehen. Einem Werk von dieser Bedeutung müssen wir eine ausführlichere Besprechung widmen.

Die Einleitung nimmt mehr als die Hälfte des ersten Bandes ein und beginnt mit der Präzession. Auf einen kurzen Bericht über ihre erste Entdeckung und die theoretischen Arbeiten von \textsc{Newton}, \textsc{d’Alembert}, \textsc{Euler}, \textsc{Lagrange} und \textsc{Laplace} folgt eine Tafel zur Veränderlichkeit der Präzession in den verschiedenen Jahrhunderten. Daran schließen sich Formeln und Hilfstafeln für die säkulare Änderung der Breite und für denjenigen Teil der säkularen Änderung der Länge an, der von Stern zu Stern verschieden ist. Ein eigener Kunstgriff des Verfassers macht diese Hilfstafeln leicht benutzbar. Es folgen die Formeln für die Präzession in Rektaszension und Deklination sowie eine Tafel für die Werte der Konstanten von 1450 bis 1950. Die Grundwerte hat der Verfasser durch eigene, sehr sorgfältige Vergleiche ermittelt. Zu Recht hat er dazu nur die neueren Beobachtungen herangezogen. Denn die Beobachtungen von \textsc{Hipparch} und \textsc{Tycho} können trotz des größeren Zeitabstands wegen ihrer weit geringeren Güte bei Weitem keine so zuverlässigen Ergebnisse liefern wie die Beobachtungen, die \textsc{Mayer} und \textsc{Bradley} mit den besten Mauerquadranten angestellt haben. Selbst \textsc{Flamsteed}s Beobachtungen bleiben noch zu weit hinter diesen zurück. Sie müssten in jedem Fall ganz neu reduziert werden, wenn man etwas Brauchbares aus ihnen folgern wollte. Der Rezensent hätte allerdings den Umstand, dass die möglicherweise noch bestehende geringe Unsicherheit bei der Abnahme der Schiefe der Ekliptik und bei der Eigenbewegung der Fixsterne bei zeitlich weiter auseinanderliegenden Beobachtungen größer sei als bei näher beieinanderliegenden, unter den Gründen für den Ausschluss der Beobachtungen von \textsc{Flamsteed} nicht angeführt. Aus dem Vergleich von 175 Deklinationen von \textsc{Bradley} mit denen von \textsc{Barry} leitet der Verfasser die gesamte jährliche Präzession von $50''1670$ ab. Der Vergleich der Deklinationen von \textsc{Mayer} und \textsc{Barry} ergibt $49''9140$. Ferner ergaben 350 Rektaszensionen von \textsc{Mayer} im Vergleich mit der eigenen Bestimmung des Verfassers $50''0718$ und schließlich 222 Rektaszensionen von \textsc{Bradley} im Vergleich mit denen des Verfassers $50''0631$. Das Mittel aller Bestimmungen beträgt somit $50''0540$, und dieser Wert liegt allen Formeln und Tafeln des Werks zugrunde. Wir erwähnen die folgenden Untersuchungen über die Eigenbewegungen der Sterne nur kurz. Man wird dort unter anderem mit Vergnügen eine ganz neue Reduktion aller Beobachtungen der Deklination des Polarsterns von \textsc{Flammsteed} finden, aus der sich nur eine säkulare Zunahme von $3''$ ergibt. Zur Berechnung der Deklination aus Länge und Breite benutzt der Verfasser hier Tafeln mit 10 Dezimalstellen. Das ist freilich nötig, wenn man die Deklination über ihren Sinus bestimmt. Diese Unbequemlichkeit lässt sich aber leicht vermeiden, und mit den gewöhnlichen Tafeln kommt man ebenso weit. Eine Tafel für die mittlere Rektaszension und Deklination des Polarsterns für jedes einzelne Jahr von 1790 bis 1820 schließt diese gehaltvolle Abhandlung ab.

\begingroup\rightskip=0pt plus 1em
Um die Konstante der \emph{Aberration} zu bestimmen, hat der Verfasser sämtliche Original%
\SourceLeaf{496}%
beobachtungen von \textsc{Bradley} erneut mit größter Genauigkeit untersucht. Er erhält $20''232$, was nur ganz geringfügig von dem Ergebnis abweicht, das \textsc{Delambre} aus den Verfinsterungen der Jupitermonde bestimmt hat, nämlich $20''255$. Den letzteren Wert hat der Verfasser einer allgemeinen Tafel für die Aberration am Ende des zweiten Bandes zugrunde gelegt. Diese Tafel nimmt nur zwei Seiten ein, aber ihre Benutzung erfordert eine Multiplikation mit einem Sinus und einem Sekans. Auch dieser kleinen Unbequemlichkeit hat der Verfasser durch einen eigenen Kunstgriff abgeholfen, sodass man bei einer anderen, in die Einleitung aufgenommenen allgemeinen Tafel nur addieren muss. Darauf folgen noch Formeln und Tafeln für die Aberration unter Berücksichtigung der elliptischen Bewegung der Erde sowie für die Aberration aufgrund der täglichen Bewegung. Schließlich folgen Formeln und Tafeln für die Änderung der Aberration bei einer Änderung der Rektaszension und Deklination des Sterns um $1^\circ$. Damit lässt sich abschätzen, wie lange die speziellen Tafeln brauchbar bleiben.\par\endgroup

Auch für die \emph{Nutation} hat der Verfasser außer den Formeln zwei allgemeine Tafeln geliefert. Die eine steht am Ende des zweiten Bandes und setzt die große und die kleine Halbachse der Nutationsellipse nach \textsc{Laplace} mit $10''056$ und $7''486$ an. Ihre Benutzung erfordert eine Multiplikation mit einem Tangens. Diese Unbequemlichkeit entfällt bei der anderen Tafel, die in die Einleitung aufgenommen wurde. Dort hat Herr \textsc{von Zach} die Halbachsen allerdings mit $9''648$ und $7''182$ angesetzt, weil er glaubte, die von \textsc{Laplace} zugrunde gelegte Mondmasse etwas verringern zu müssen. Außerdem findet man in diesem Abschnitt Anleitungen zur Berechnung des von der Sonnenlänge abhängigen Teils der Nutation sowie Formeln und Tafeln für die Änderung der Nutation bei veränderter Position des Sterns.

Der übrige Teil der Einleitung enthält Anleitungen zur richtigen Einstellung des Passageninstruments, Erklärungen zu den Verzeichnissen und Tafeln in diesem Band sowie Untersuchungen und Tafeln zur Aberration der Planeten, zu ihrer Parallaxe und zu ihren Durchmessern. Der Behauptung auf S. 203 können wir nicht zustimmen, wonach das Verfahren zur Berechnung der Aberration der Planeten, das wir allen anderen vorziehen und bei dem man nur die Zeit ändert, --- lediglich den von der Bewegung des Planeten abhängigen Teil der Aberration liefert. Dabei versteht sich, dass auch die verwendete Position der Erde dem geänderten Zeitpunkt entsprechen muss.

\begingroup\rightskip=0pt plus 1em
Übrigens sind alle Anleitungen in dieser gehaltreichen Einleitung durch zahlreiche gut gewählte Beispiele so erläutert, dass selbst die Ungeübtesten bei ihrer Anwendung auf keine Schwierigkeiten stoßen werden. Überall sind die Schriften angegeben, in denen man sich über die behandelten Gegenstände weiter unterrichten kann.\par\endgroup

Von den nun folgenden Tafeln selbst geben wir den Inhalt nur im Überblick wieder. Sie enthalten die Rektaszensionen und Deklinationen der 36 wichtigsten Sterne nach der neuesten Bestimmung von \textsc{Maskelyne} sowie spezielle Aberrations- und Nutationstafeln für jeden dieser Sterne. Es folgen ähnliche Tafeln für den Polarstern für 1790, 1800, 1810 und 1820 sowie ein Verzeichnis von Fixsternen, deren Rektaszensionen fast gleich und deren Deklinationen fast gleich groß, aber entgegengesetzt sind. Dieses Verzeichnis dient dazu, die Stellung des Passageninstruments zu berichtigen. Darauf folgt der wertvollste Teil des ganzen Werks, nämlich die an der Seeberger Sternwarte bestimmten Rektaszensionen von 1830 Tierkreissternen für das Jahr 1800 mit den Abweichungen von \textsc{Piazzi}s Katalog. Diese Abweichungen sind meist gering und werden in der Regel noch viel geringer, wenn man zu \textsc{Piazzi}s Rektaszensionen die Korrektur von $3''8$ addiert. Dennoch ist es stets gerechtfertigt, den Bestimmungen den Vorzug zu geben, die mit dem prächtigen 8 Fuß langen Seeberger Meridianfernrohr gemacht wurden. Die Deklinationen sind hier nur zur Identifizierung der Sterne angegeben, und zwar lediglich in Minuten. Bei einigen Sternen fehlen sie ganz. Es wäre zu wünschen, dass sich jemand die Mühe machte, diese Lücke zu schließen. Meist würde dafür schon die Histoire céleste genügen, sofern nicht zufällig mehrere Sterne im Tierkreis bis auf eine Zeitsekunde in der Rekta%
\SourceLeaf{497}%
szension übereinstimmen. Viele dieser Sterne scheinen sogar in dem unmittelbar folgenden Deklinationsverzeichnis von \textsc{Barry} vorzukommen, zum Beispiel in der Jungfrau 921 und 963. Um auch eine hinreichende Anzahl genau bestimmter Deklinationen zu geben, hat Herr \textsc{von Zach} ein zweites Verzeichnis von etwa 1200 Tierkreissternen beigefügt. Es enthält nur diese Deklinationen und die ungefähren Rektaszensionen und gibt das Ergebnis der Beobachtungen wieder, die \textsc{Barry} und \textsc{Henry} an der Mannheimer Sternwarte mit einem 8 Fuß langen Mauerquadranten von \textsc{Bird} gemacht haben. Die ebenfalls angegebenen Abweichungen von \textsc{Piazzi}s Katalog zeigen, dass dieses Verzeichnis es durchaus verdient, ihm zur Seite gestellt zu werden. Auch hier fehlen bei mehreren Sternen die Rektaszensionen ganz. Das fällt hier aber weniger ins Gewicht, weil meist schon die Reihenfolge zur Identifizierung des Sterns dienen kann.

Der verbleibende Teil des ersten Bandes enthält eine allgemeine Tafel für die Präzession in Rektaszension und Tafeln für die mittleren Refraktionen und deren Korrektur für den Stand des Barometers und Thermometers. Die Refraktionstafeln folgen denselben Grundsätzen wie die Tafeln in \textsc{Delambre}s neuen Sonnentafeln und sind lediglich auf die bisher allgemein üblichen Maße umgerechnet. Dazu kommen Tafeln für die Sonnenparallaxe und die Umrechnung der Sternzeit in mittlere und wahre Sonnenzeit sowie einige weitere kleinere Tafeln.

Den \emph{zweiten Band} füllen fast ganz die speziellen Aberrations- und Nutationstafeln für 494 Tierkreissterne samt ihren Positionen für 1800 nach dem Verfasser und anderen Astronomen. Wie wertvoll dieses Hilfsmittel für praktische Astronomen ist, haben wir schon erwähnt. Es bleibt also nur anzumerken, dass der zweite Band schon vor dem ersten gedruckt worden ist und daher in verschiedenen Punkten von ihm abweicht. Die vom Verfasser stammenden Sternpositionen müssen daher sämtlich durch die im großen Katalog des Bandes enthaltenen Positionen ersetzt werden, die nach \textsc{Maskelyne}s Verbesserung seines Fundamentalverzeichnisses berichtigt sind. Sollen ferner die Präzessionen, Aberrationen und Nutationen mit den neuesten Bestimmungen im ersten Band vollständig in Übereinstimmung gebracht werden, so sind die ersten um $\frac{1}{124}$ zu vermindern, die zweiten um $\frac{1}{78}$ zu erhöhen und die dritten um $\frac{1}{92}$ zu erhöhen.

Den Abschluss des zweiten Bandes bilden die bereits erwähnten allgemeinen Tafeln für Aberration und Nutation sowie eine Tafel für die Länge des aufsteigenden Knotens der Mondbahn nach den Mondtafeln von \textsc{Bürg}.
